\documentclass[12pt]{article}
\usepackage{graphicx} % Required for inserting images
\usepackage{float} % For [H] float placement
\usepackage{booktabs}
\usepackage{caption}
\usepackage{adjustbox}
\usepackage[utf8]{inputenc}  % This is for old compilers, but good practice
\usepackage[T1]{fontenc}    % This is for correct character encoding (accents, ñ, etc.)
\usepackage[margin=3cm]{geometry} % This line changes the margins
\usepackage{subcaption}
\usepackage{url}
\usepackage[most]{tcolorbox} % paquete para cuadros
\usepackage{xcolor}
\usepackage{setspace}
\usepackage{amssymb}

%for citations
\usepackage{csquotes}       % Requerido por biblatex para manejar comillas según el estilo
\usepackage[style=apa, backend=biber]{biblatex}
\addbibresource{references.bib}
\usepackage{hyperref}       % Siempre cárgalo DESPUÉS de biblatex

%for notes and corrections
\usepackage{soul}
\usepackage{xcolor}
\usepackage{comment}
%para tablas largas
\usepackage{longtable}

\onehalfspacing
\title{Cuando OXXO llega a la cuadra: }
\author{Sophia Aristizábal Flórez}
\date{\small{\url{https://github.com/sophiearisti/proyecto_econometria_avanzada_real.git}}}


\begin{document}

\maketitle
\onehalfspacing

\section{Abstract}

\newpage

\section{Acknowledgements}

\newpage

\section{Introduction}

The rapid expansion of modern retail chains in Latin America after the COVID-19 pandemic has transformed urban commercial dynamics and people's daily routines in the main cities of the region \parencite{Antom2024,McKinsey2024,MDPI2021}. Convenience stores’ growing commercial importance has prompted economic research to primarily focus on the impact of these establishments on retail businesses, consumer behavior, and the labor market \parencite{delgado2024}. However, beyond their economic role, their physical and operational characteristics are associated with mixed findings in the literature on urban crime.\\ 

These establishments are strategically located in high-traffic areas, and commercial corridors. Additionally, these 24-hour stores offer basic goods, financial services, and bill payment services, attracting diverse customer flows throughout the day and night. Given their expanding footprint within the urban landscape, these retail hubs have naturally emerged as magnets of economic, urban and pedestrian activity \parencite{Dugato2026,marcos2022}.\\

Given these characteristics and depending on the context, these stores are seen as \textit{crime generators} because they increase foot traffic and create incidental criminal opportunities \parencite{Dugato2026}. Also, \cite{Wellford1997} identify them as robbery-prone commercial establishments because they have available cash on hand. Nonetheless, \cite{Davis2021} offer a different perspective, showing that convenience stores are seen as crime deterrents by providing natural surveillance and increasing “eyes on the street” through extended operating hours and improved lighting. This duality highlights the necessity of analysing their role in the surrounding criminal environment, as targeting the aspects that facilitate criminal behaviour is often more effective for long-term crime reduction than traditional repression \parencite{Dugato2026}.\\

Within this retail landscape, OXXO, a Mexican chain now owned by FEMSA, represents a valuable case study of of how the rapid proliferation of convenience stores interacts with urban  criminal activity. While OXXO maintains a massive regional presence of over 25,000 locations across Latin America, its footprint in Colombia is strikingly centralized, with over 300 of the more than 500 stores in the country concentrated in the capital, Bogotá \parencite{godoy2025}.This high density within a single metropolitan area provides a unique "urban laboratory" to conduct a granular, interurban crime analysis. By exploiting the staggered arrival of these establishments across Bogotá’s Zonal Planning Units (UPZs), this thesis evaluates the causal impact of OXXO openings on local criminal activity. Specifically, I examine how these effects vary by crime category, and whether the impact intensifies or dissipates over time following the store's inauguration.

To estimate this impact, I used Colombia’s public georeferenced crime data stored in the Statistical, Criminal, Contraventional, and Operational Information System of the National Police (SIEDCO). I created a quarterly panel between 2015 to 2022 aggregated at the level of the UPZ, which serves as a sub-district administrative division in Bogotá. I also employed Two-Way Fixed Effects (TWFE) and the Callaway and Sant'Anna (C\&S) estimator, along with their corresponding event study methods, to capture the dynamic evolution of the treatment effect over time.\\

With this empirical strategy, I found that there is no effect of staggered OXXO entry on crime activity across the city. Using both aggregated (TWFE) and dynamic specifications, I showed that OXXO had no impact on different types of crime across Bogotá. Additionally, after considering other similar establishments located near OXXO’s stores, the effect remained insignificant. These results may suggest that, in the long and short run, OXXO or convenience stores do not affect Bogotá’s security dynamics. This may be because they enter areas that are already highly frequented throughout the day, which makes OXXO an insignificant factor in influencing crime rates in the surrounding areas.\\

In general, researching this phenomenon is crucial for the city. Although Bogotá has experienced significant crime reduction since the 1990s \hl{()}, it is still a major public policy concern as high volumes of personal theft and a rising public perception of insecurity continue to affect the daily lives of its citizens \parencite{CCB2026}. Additionally, crime remains spatially concentrated, with substantial variation across zones driven by socioeconomic stratification, land use, and police presence \parencite{Enriquez2014,GimenezSantana2018}, wich underscores the importance of evaluating how the introduction of these modern retail chains could reconfigure Bogotá’s urban criminal environment.\\

Despite extensive research on urban crime in Bogotá, no study has systematically examined the causal impact of convenience store proliferation, particularly OXXO, on local crime patterns. Existing research has focused on public transport infrastructure \parencite{Bacares2013a,Bacares2013b,Garcia2005}, large-scale police interventions \parencite{Blattman2017,Blattman2021}, alcohol sales restrictions \parencite{Mello2013}, and socioeconomic determinants \parencite{Escobar2012}.One recent study examined the relationship between hard-discount stores (Tiendas D1) and crime \parencite{Lopez2024}, but OXXO represents a distinct retail model characterized by 24-hour operations and a higher density of ``grab-and-go'' items, which may create different criminal opportunities compared to the limited operating hours of hard-discount formats.\\

In the next section I briefly introduce OXXO’s expansion across Bogotá, and the city’s crime rates over the years. Then, in Section 3 I present the related literature about the spatial nexus of crime and convenience stores. In Sections 4 and 5 I describe the data sources used and the general descriptive statistics. After that I explain in Section 6 the econometric models used and in Section 7 I show the results obtained. In Section 8 I include some robustness checks related with the unit of analysis used, the control and treatment groups, and the specification of some variables. Finally, Section 9 discusses probable mechanisms that may influence the results obtained.

\section{Background}
\subsection{The expansion of OXXO in Bogotá}
OXXO is one of the most prominent convenience store brands in Latin America, and since its arrival to Colombia in 2009 it has consolidated itself as one of the main convenience store chains in the country. 
Although OXXO has been operating in Colombia for almost 17 years, it only gained momentum after the pandemic. 
In fact, in Bogota the number of establishments skyrocketed from 82 in 2020 to 384 in 2025. 
This represents an increment of over 368\% in only five years. Figure~\ref{fig:EVOLUTION_OXXO} depicts its accelerated growth in Bogotá, as it shows the number of OXXO locations in the capital by year since its arrival.\\

\begin{figure}
    \centering
    \includegraphics[width=0.9\linewidth]{other_figs/oxxos_por_ano.png} % Ajusta el path y tamaño
    \caption{
        \textbf{Yearly evolution of the number of OXXOs in Bogotá} 
        \newline
     }   
    \label{fig:EVOLUTION_OXXO}
\end{figure}

OXXO’s main objective is to be part of people’s lifestyles by offering an ample variety of products and services that can be useful and practical for their daily lives \parencite{oxxo_quienes_somos}. 
Another distinctive characteristic is that 95\% of its establishments function in a 24-hour format \parencite{oxxo_faqs}, providing a competitive advantage by being the primary option of consumers whenever they need to satisfy spontaneous consumption needs, regardless of the hour. 
These well-defined competitive objectives and inherent brand characteristics have allowed OXXO to proliferate rapidly throughout Bogotá, where it has been widely embraced by the local population.\\

In line with this strategy, this modern retail chain began its expansion in the northern part of the city, later expanding throughout the eastern region. 
This strategic behavior responds to the inherent characteristics of the northeastern neighborhoods of Bogotá, where a higher concentration of formal employment and high-income residential clusters provide an ideal environment for the convenience store model  \hl{(citation)}.\\

Figure~\ref{fig:fourpanelOXXO_evolution_UPZ} presents this evolution, and shows that now, OXXO is present in almost 70\% of all the city´s Zonal Planning Units (UPZ). 
UPZs are analysis units that group neighborhoods based on their socioeconomic stratification, commercial, land use, and housing characteristics, and living conditions of their inhabitants \parencite{alcaldia_bogota_upz}. 
Currently, the only areas without OXXO establishments are in the far southern reaches of Bogotá, which correspods to UPZs with low levels of economic performance, socioeconomic stratification, urban transport connection, and are blocks that continue to relly heavily in traditional neighborhood stores \hl{(citation)}.\\

\begin{figure}
    \centering
    % Top row
    \begin{subfigure}[b]{0.3\textwidth} % Reduced from 0.4 to 0.3
        \includegraphics[width=\linewidth]{figs_oxxo_maps/mapa_oxxos_binary_upz_2009.png}
        \caption{2009}
        \label{fig:panelA1}
    \end{subfigure}
    \hspace{0.5cm} % Adds specific horizontal space between images
    \begin{subfigure}[b]{0.3\textwidth} 
        \includegraphics[width=\linewidth]{figs_oxxo_maps/mapa_oxxos_binary_upz_2011.png}
        \caption{2011}
        \label{fig:panelB1}
    \end{subfigure}
    \hspace{0.5cm} % Adds specific horizontal space between images
        \begin{subfigure}[b]{0.3\textwidth} 
        \includegraphics[width=\linewidth]{figs_oxxo_maps/mapa_oxxos_binary_upz_2015.png}
        \caption{2015}
        \label{fig:panelC1}
    \end{subfigure}
    
    \vspace{0.3cm} % Vertical space between rows
    
    % Bottom row
    \begin{subfigure}[b]{0.3\textwidth} 
        \includegraphics[width=\linewidth]{figs_oxxo_maps/mapa_oxxos_binary_upz_2019.png}
        \caption{2019}
        \label{fig:panelD1}
    \end{subfigure}
    \hspace{0.5cm}
    \begin{subfigure}[b]{0.3\textwidth} 
        \includegraphics[width=\linewidth]{figs_oxxo_maps/mapa_oxxos_binary_upz_2023.png}
        \caption{2023}
        \label{fig:panelE1}
    \end{subfigure}
    \hspace{0.5cm}
    \begin{subfigure}[b]{0.3\textwidth} 
        \includegraphics[width=\linewidth]{figs_oxxo_maps/mapa_oxxos_binary_upz_2025.png}
        \caption{2025}
        \label{fig:panelF1}
    \end{subfigure}

    \vspace{0.3cm} % Vertical space between rows

    \begin{subfigure}[b]{0.3\textwidth} 
        
\includegraphics[width=\linewidth]{figs_oxxo_maps/legend.png}
        \label{fig:legend}
    \end{subfigure}
    
    \caption{\textbf{Evolution of the presence of OXXO by UPZ}}
    \footnotesize{\textit{Note:} This map only indicates which UPZs have OXXOs, it does not consider the number of locations present in each Unit}
    \label{fig:fourpanelOXXO_evolution_UPZ}
\end{figure}

As noted, OXXO’s stores are not randomly located. 
They are centrally administered by FEMSA (meaning the brand does not operate under a franchise model) with the purpose of standardizing processes, logistics, services, and products \parencite{oxxo_faqs}. 
This means that each location is carefully planned to respond to the overall interests of the company. In fact, if we attentively analyze how these establishments are clustered around Bogotá, we can identify that they tend to localize in areas where there is higher traffic and spatial accessibility. This helps OXXO to maximize its market reach and ensure high levels of consumer visibility.\\

Figure~\ref{fig:oxxo_transmilenio_vias} graphically demonstrates this correlation by plotting the city’s arterial roads, main avenues, and Transmilenio stations (Bogotá's Bus Rapid Transit system); 
it shows that by 2025, OXXO locations tend to group along these major transport corridors.\\

\begin{figure} [H]
    \centering
    \includegraphics[width=0.9\linewidth]{figs_oxxo_maps/mapa_oxxos_transmilenio_vias.png} % Ajusta el path y tamaño
    \caption{
        \textbf{Bogota's OXXO stores, Transmilenio stations and main roads in 2025} 
        \newline
     }   
    \label{fig:oxxo_transmilenio_vias}
\end{figure}

With these in mind, it is reasonable to think there may be a correlation between urban crime and modern convenience stores chains, especially in the case of OXXO. 
Locating in areas with high foot traffic means more potential targets for crimes such as robbery or assault. 
Also, due to their 24-hour format and alcohol sales during the night, OXXO locations serve as  informal gathering points that may escalate into conflicts or provide opportunities for illicit activities. 
Additionally, locations near major transit areas and public transport stops increase accessibility for both legitimate customers and offenders \parencite{Dugato2026, shin2024dollar, feng2019aggravating}.

\subsection{Criminality in Bogotá}

Between the 80s and 90s Bogota was recognized as one of the most insecure cities of Colombia, in fact, in 1991 it reached an annual crime rate of 1,287 crimes for each 100 thousand inhabitants \parencite{rodriguez2017ciudad}. 
During this time the capital lived an upswing of narcoterrorism and the assassination of political leaders. However, at the end of the 90s and beginning of the 2000s, crime started to decrease; for example, between 1993 and 2002 Bogotá saw an estimated 65\% decrease in the homicide rate, while robbery fell by about 33\%. 
This shift was driven by citizen culture programs, large-scale urban renewal projects, and robust institutional reforms implemented by the majors of this period \parencite{Sanchez2003, mockus2012building}. 
Finally, in past years urban crime has shown resistance to keep declining. After the 2019 covid pandemic criminality had a significant rebound, with personal theft surpassing pre-pandemic levels. By 2025, despite institutional efforts, the city's crime rates have shown resistance to reduce, particularly regarding homicides and robberies hl{(citation)}.\\

Currently, Bogota’s crime, as any other city, clusters in specific zones. 
The most insecure Localities (a unit of territorial division that gathers many UPZ) are located in the southeastern region of Bogota, wich correlates with a high presence of criminal organizations, informality, and deficencies of human development \parencite{probogota2024indice}. 
However, this concentration does not imply that all crime categories peak in the same areas. 
While homicide rates tend to follow this general rule in the city, personal theft exhibits a distinct geographic distribution, primarily clustering in the city's northeastern sector \parencite{probogota2022informe}.\\ 

Figure~\ref{fig:oxxo_transmilenio_vias} shows this behaviour as it displays how personal theft rates distribute throughout Bogotá´s UPZs over the years. 
These maps also highlight that between 2015 and 2021 (and even later, as it can be noted in ProBogota’s annual security informs hl{(PIE DE PAGINA)}) this tendency has not drastically changed, meaning that there are environmental characteristics that keep favoring robberies or attracting offenders. 
Other characteristics that these zones have are high pedestrian throughput, the density of commercial land use, and the proximity to the city's main transit arteries.\\ 

\begin{figure}
    \centering
    % Top row
    \begin{subfigure}[b]{0.3\textwidth} 
        \includegraphics[width=\linewidth]{figs_dep_maps/mapa_upz_gpkg_all_years_2015_theft_to_people_index_eb_edited.png}
        \caption{2015}
        \label{fig:panelA2}
    \end{subfigure}
    \hspace{0.5cm} 
    \begin{subfigure}[b]{0.3\textwidth} 
        \includegraphics[width=\linewidth]{figs_dep_maps/mapa_upz_gpkg_all_years_2017_theft_to_people_index_eb_edited.png}
        \caption{2017}
        \label{fig:panelB2}
    \end{subfigure} % <--- THIS WAS MISSING
    
    \vspace{0.3cm} 
    
    % Bottom row
    \begin{subfigure}[b]{0.3\textwidth} 
        \includegraphics[width=\linewidth]{figs_dep_maps/mapa_upz_gpkg_all_years_2019_theft_to_people_index_eb_edited.png}
        \caption{2019}
        \label{fig:panelC2} % Fixed label sequence (A, B, C, D)
    \end{subfigure}
    \hspace{0.5cm}
    \begin{subfigure}[b]{0.3\textwidth} 
        \includegraphics[width=\linewidth]{figs_dep_maps/mapa_upz_gpkg_all_years_2021_theft_to_people_index_eb_edited.png}
        \caption{2021}
        \label{fig:panelD2}
    \end{subfigure}
    
    \caption{\textbf{Evolution of theft rates by UPZ}}
    \label{fig:fourpanelCrime}
    \footnotesize{\textit{Note:} The theft rate per 10,000 inhabitants for each unit was calculated using rate smoothing techniques (Empirical Bayes smoothing).
    This approach aimed to reduce the influence of outlier UPZs that had remarkably high rates due to their small populations rather than high theft activity.
    Still,after this transformation, I had to exclude two UPZs and set their rates to zero, as they continued to show extreme values that overshadowed the crime activity of other areas. 
    In fact, one UPZ is an airport, and the other is mainly rural and contains many storage facilities instead of dwellings.}
\end{figure}

Factors that favor crime clustering must be taken into account because if OXXO stores also locate in inherent high crime areas, it will be important to evaluate whether OXXO efectively causes an increment in crime or if both OXXO and criminals are simply attracted to the same high-activity "nodes."

\section{Related Literature}

In the context of urban economics, spatial econometrics and crime literature, various studies have explored how built environment features and infrastructure shape crime activity in cities. A number of empirical analyses have investigated this by exploring the causal impact of urban public transit on crime. 
\cite{ijgi10100632}, for example, assessed the impact of two new subway lines implemented in 2017 in ZG city, China, on thefts in the surrounding areas. 
By employing a Difference-in-Differences (DiD) framework, they found that in the short run, theft increased by 15\%.
Others, on the other hand, have ventured to explore the paper of new urban policies and infrastructure interventions on this topic. 
\cite{dominguez2022crimetime} paper is one case, as they study the causal effect of ambient light on crime in Santiago, Chile, by using the "Daylight Saving Time" policy as a natural experiment. Due to this policy, they could find that an increase of one hour of ambient light exposure between 7 and 9 reduced thefts by 30\%. 
\hl{These studies empahsize on XXXXX (give emphasis on a crime literature concept)}\\ 

Retail chains can be considered as built environment features that may influence urban crime activity. 
While different studies have extensively explored their role, these previous inquiries have their own limitations. 
A substantial body of research dates back to the pre-2000s era, meaning that their findings do not contemplate the current urban dynamics of modern cities.
Nonetheless, these foundational studies provide a critical perspective on the underlying mechanics of crime. For instance, early investigations focused on identifying the characteristics that influenced the propensity of convenience stores to be targeted for robbery. 
They identified that stores that operated for longer hours at night were more likely to be robbed, as there are fewer people around \parencite{dalessio1990crime}. 
Others found that these types of stores were more prone to be robbed due to the lack of visible security such as police presence, security guards, or effective alarm systems \parencite{Wellford1997, Dugato2026}. 
Although useful, they fail to account for convenience stores’ role on crime of their immediate surroundings.\\

On the other hand, other recent publications have examined the influence of analogous commercial establishments on urban crime. 
While some of these studies do not focus exclusively on convenience stores, they utilize other retail outlets with similar operational profiles which can serve as proxies for what this thesis aims to investigate. 
\textcite{askey2018fast} is and example as the authors pooled data from fast-food restaurants and convenience stores in Seattle. They used sales revenue as a proxy for human activity to count the different levels of operational intensity across locations, and aimed to explain crime patterns in the immediate surroundings of these establishments. 
Their results indicate that stores with higher average sales were associated with with a 7.4\% increase in expected crime counts. 
Conversely, \cite{feng2019aggravating} contribute to this discussion by finding that alcohol outlets influenced the likelihood of aggravated assault and street robbery in New York City. 
Their findings reveal differential impacts depending on the type of outlet (if it is a bar, a grocery store, or other similar establishment); for instance, the authors highlighted that in Manhattan and the Bronx, proximity within two blocks of a grocery store increased the relative risk of street robbery by more than sixfold.\\ 

Another relevant investigation is the one of \parencite{shin2024dollar}. 
Shin used a DiD approach to invetigate the impact of dollar stores on localized violent crime patterns in Chicago. 
To ensure the treatment and control groups were as homogeneous as possible, the author created localized areas using a 250-feet buffer for the treatment group and a 250-to-354-feet outer ring as the control. 
At the end, this research found that store openings lead to a 21.\% increase in violent crime (0.17 incidents per quarter), a result primarily driven by a 38\% surge in robberies.\\

Although these studies offer a useful baseline of how convenience stores, and specifically OXXO, impact nearby crime, they present an incomplete perspective as they exclusively explored this nexus in cities of high-income nations. 
To build a more robust theory, these hypotheses must be tested within Latin American cities which have particular urban dynamics. Many papers argue that convenience stores are normally located in socio-economically disadvantaged areas where poverty, social disorganization, and lack of community resources exacerbate crime risks \parencite{shin2024dollar, Dugato2026}, but in cities like Bogotá the situation is not analogous. 
OXXO in Bogotá, for instance, does not follow this rule as it instead prioritizes to open locations in middle to high income neighborhoods, and places with high commercial activity and influx of people.\\ 

In Latin America, the economic literature regarding convenience stores remains limited. Recent investigations have explored their causal role on various labor market indicators \parencite{delgado2024, marcos2022}. 
However, despite this growing interest in labor dynamics, the spatial relationship between these establishments and urban security in the region has been largely overlooked. Specifically in Bogotá crime literature has primarily focused on public transport, general socioeconomic urban caracteristics, and policy intervention impact on crime \parencite{moreno2005impacto, bacares2013impact}.\\ 

The investigations that have actually explored this topic are counted. 
For example,  \cite{jaramillo2021low} used the same method as \parencite{shin2024dollar} to study the effect of D1 hard discount stores on the crime around the streets where they open in Bogotá. 
Although D1s are not convenience stores per se, they attract a similar type of costumers and and often share similar spatial locations with OXXO outlets. By analyzing the period of 2016 to 2019 and employing the same information utilized in this thesis (SIEDCO), the authors found that D1 openings caused a 33.4\% average reduction in crime density per block. 
They attributed this impact to the 'eyes on the street' phenomenon, suggesting that the increased pedestrian flow acts as a form of informal guardianship that deters criminal activity \parencite{ jacobs1961death, jaramillo2021low}.\\ 

\section{Data used and unit of analysis}

For this thesis I use UPZs as the unit of analysis. UPZs work as a feasible division of Bogotá as the disctrict maintains continuity and consistency of data across the years, and security interventions are normally focalized by UPZ \hl{(cita)}. 
Also, these Units maintain urbanistic and socioeconomic homogeneity by grouping similar neighborhoods, and are the middle point between Localities and Street blocks (manzanas).\\ 

Meanwhile, the results of this study employ four main data sources (see appendix~\ref{sec:variables_descripcion_detallada}).
For measuring crime rate in Bogotá I used the georeferenced crime events stored in SIEDCO \footnote{Access at \url{ https://oaiee.scj.gov.co/agc/rest/services/Tematicos_Pub/SIEDCO_Delitos_Pub/FeatureServer}} spanning from \hl{2018 to 2022}. 
SIEDCO is Colombia’s official Police database that consolidate real-time citizen reports and administrative records of criminal activity. Here, for each incident recorded the system saves details including type of crime (e.g., homicide, personal robbery, commercial burglary, or cell phone theft), the precise geographic coordinates of the event, date and time of the incident, and the victim’s gender. 
With this information, I aggregated the crime data at the UPZ-year level, categorizing incidents by offense type. To ensure comparability across years and Units, I normalized these counts using annual population estimates for each UPZ. This allowed for the construction of a crime rate per 10,000 inhabitants, as shown in this equation:\\

\begin{equation}
crime\_type\_index_{upz,year}=\frac{total\_incidents_{upz,year}}{population_{upz,year}} ; \quad \forall upz,year
\end{equation}

The period 2018 to 2022 is adequate for the analysis as they englobe the years in which OXXO had its explosive expansion across Bogota. 
This timeframe ensures sufficient temporal and spatial variability, meaing it allows for a higher proportion of UPZs to enter treatment during the study period compared to those always treated. 
Prior to 2018, the lack of new store openings in other UPZ would result in an insufficient number of newly treated units, limiting the identifying variation necessary for using TWFE.\\ 

Additionally, from the Registro Único Empresarial y Social (RUES) \footnote{Access at \url{ https://www.rues.org.co}}, the central business information system in Colombia that contains all the records of companies and merchants registered by the chambers of commerce in the country, I identified the opening and closing year of each OXXO. 
Additionally, with Google Maps API I obtained the precise georeferenced location of each store. Concerns may arise on how precise was this match of each store with its correct location, specially for the ones that no lorgers exist. 
I manually checked the validity of each match with official OXXO publications, and as a result \hl{XX\%} of establishments were properly matched.\\ 

I also did this same process with D1, Justo\&Bueno and ARA hard discount chains. The reason why I also took into account these brands is because they often share similar spatial locations with OXXO outlets; in fact, retailers tend to cluster in certain areas to take advantage of agglomeration benefits \parencite{Konishi2005, Seong2022}. 
Given that prior research suggests these establishments may influence both crime and commercial activity \parencite{delgado2024, jaramillo2021low}, their inclusion as control variables is important to evaluate the validity of the results.\\

\hl{Finally, I used the Quality of Life Survey of Bogotá 2007 \footnote{Access at \url{https://www.dane.gov.co/index.php/estadisticas-por-tema/salud/calidad-de-vida-ecv/encuesta-nacional-de-calidad-de-vida-2007-bogota}} to analyze the initial socioeconomic characteristcs of never treated and treated UPZ before OXXO arrived in 2009. 
And with Mapas-Bogotá I obtained the fixed variable socioeocnomic and environmental controls at the UPZ level just to explore how are OXXOs not randomly located, as these will be inhenrently controlled by TWFE.}

\section{Descriptive Statistics} 


\section{Control Variables}

\begin{table} [H]
  \centering
  \caption{ Differences of means between treated and never treated for fixed controls}
  \label{tab:diferencias_medias_controles_fijos}
  \begin{adjustbox}{width=\textwidth,center}
    \begin{tabular}{l c c c}
      \toprule
      \multicolumn{1}{c}{} & \multicolumn{1}{c}{\textbf{Never treated}} & \multicolumn{1}{c}{\textbf{Treated}} & \multicolumn{1}{c}{\textbf{Difference of means}} \\
      \midrule
       \midrule
      \multicolumn{4}{l}{\textbf{Baseline}} \\
      \midrule
      Inhabitants by locality in 2007 & 503.534 & 482.463 & -21.071 \\
      & (38.331) & (43.823) & \\
      Household size 2007 & 3.714 & 3.335 & $-0.380^{***}$ \\
      & (0.049) & (0.048) & \\
      Index of quality of life 2007 & 89.121 & 91.724 & $2.603^{***}$ \\
      & (0.151) & (0.216) & \\
      Average monthly expenditure 2007 & $8.97e+05$ & $1.38e+06$ & $4.82e+05^{***}$ \\
      & $(55836.359)$ & $(57218.217)$ & \\
      \midrule
      \multicolumn{4}{l}{\textbf{Fixed controls}} \\
      \midrule
      Average stratum & 2.162 & 3.352 & $1.190^{***}$ \\
      & (0.135) & (0.119) & \\
      Nearby transmilenio stations & 2.339 & 6.067 & $3.727^{***}$ \\
      & (0.404) & (0.544) & \\
      Access to arterial roads & 29.286 & 44.200 & $14.914^{***}$ \\
      & (4.490) & (2.738) & \\
      \midrule
      \midrule
      UPZs & 56 & 60 & 116 \\
      \bottomrule
    \end{tabular}
  \end{adjustbox}
   \parbox[t]{\textwidth}{%
    \vspace{0.5em}
    \footnotesize{ \textit{Note:} Significance: *** p<0.01, ** p<0.05, * p<0.1.}}
\end{table}

\section{Empirical strategy}

This research explores with two main econometric models the causal effect of the of OXXO openings on local crime activity in Bogotá. The first one I used is TWFE. 
This econometric strategy takes into account that the the arrival of OXXO’s occurs in a staggered manner. 
Additionally, it inherently controls for fixed effects for unit (UPZ) and time. 
The former is useful for capturing time-invariant unobserved characteristics of each UPZ (such as geographic location, historical infrastructure or fixed socieconomic and urban characteristics during the period of analysis), while the latter absorbs common shocks that affect all units simultaneously (such as city-wide security policies or macroeconomic fluctuations). 
Although TWFE is and adequate method for estimating the average treatement effect (ATT) of OXXO on crime activity, two identification assumptions are needed to prevent bias. 
This ones will be considered later, but let us first introduce the basic regression form of TWFE:

\begin{equation}
Y_{i,t} = \alpha_i + \alpha_t  +\beta^{twfe} D_{i,t} + \varepsilon_{i}
\label{eq:twfe}
\end{equation}

Where $ Y_{i,t} $ represents the crime rate for UPZ $i$ in quarter $t$. 
This rate corresponds to either a specific crime category or the aggregate of all crimes analyzed in this study. 
$D_{i,t}$ is a binary indicator that equals one if UPZ $i$  has at least one OXXO store in quarter $t$ , and zero otherwise. 
Furthermore, as Figure \hl{XXXX} suggests that intensity of treatment may be a factor that influences crime. 
Section \ref{sec:Results} also evaluated $D_{i,t}$ as the number of OXXO establishments that an UPZ has at a certain quarter. 
Finally, $\alpha_i$ and  $\alpha_t$ are the unit and time fixed effects respectively.\\ 

A critical identification assumption for the TWFE estimator is the homogeneity of treatment effects. 
This implies that these effects must be constant among treatment cohorts and over time \parencite{deChaisemartin2023survey}. 
This assumption may be violated if treatment effects vary based on UPZs’ characteristics or if the impact of an OXXO opening on crime evolves as the store becomes established in the area. 
Consequently, to diagnose the potential bias arising from these heterogeneities, I employ the Goodman-Bacon decomposition \parencite{GoodmanBacon2021}.  
This theorem decompoeses the $\beta^{twfe}$ as a weighted average of all possible 2×2 difference-in-differences comparisons among cohorts, and the primary concern in this decomposition is to avoid having negative weights as they indicate that the estimator may be biased from heterogeneous treatment effects \parencite{GoodmanBacon2021}.\\

Parallel trends is the other condition to guarantee an unbiased ATT, as it ensures that, in the absence of treatment, the average outcomes of the treated and control groups would have followed the same trajectory over time. 
To evaluate the validity of this assumption, I implement an Event Study design. This model requires no anticipation to treatment, and in this context it is highly unlikely that perpetrators or residents altered their behavior in expectation of an OXXO opening, as these commercial developments are typically not publicized far enough in advance to influence local criminal dynamics beforehand. 
The Events Study is estimated by using the following equation:

\begin{equation}
Y_{i,t} = \alpha_i + \alpha_t 
+ \sum_{k=-K}^{-2} \gamma^{lead}_k D^k_{i,t} 
+ \sum_{k=0}^{L} \gamma^{lag}_k D^k_{i,t} 
+ \varepsilon_{i,t}
\label{eq:eventstudy}
\end{equation}

Equation\eqref{eq:eventstudy} allows to analyze the significance of the coefficients  $\gamma^{lead}_k$  and suggests the existence of parallel trends before the arrival of OXXO. 
Additionally, the coefficients $\gamma^{lag}_k$ reflect the temporal dynamics of the impact after the treatment.\\

Finally, for robustness purposes, I will also estimate the Callaway \& Sant’Anna (CS) model. This approach accounts for treatment effect heterogeneity by estimating ATTs by cohort and period, subsequently combining these effects through a weighted average \parencite{RiosAvila_Callaway_SantAnna_2021_csdid_slides,Callaway_SantAnna_2021}. o achieve this, the model utilizes the following doubly robust estimator for the ATT of each cohort $g$ at time $t$:

\begin{equation}
ATT(g, t) = \mathbb{E} \left[ \left( \frac{G_g}{\mathbb{E}[G_g]} - \frac{\frac{p_{g,t}(X)(1 - D_t)}{1 - p_{g,t}(X)}}{\mathbb{E} \left[ \frac{p_{g,t}(X)(1 - D_t)}{1 - p_{g,t}(X)} \right]} \right) (Y_t - Y_{g-1}) \right]
\end{equation}

Where $t$ is a particular time period, $ G_g$ is a binary indicator for whether an observation belongs to cohort $g$, $ D_t $ treatment status at time $t$, and  $(1 - D_t)$ identifies observations that have not yet been treated, serving as the control group. 
Additionally, the approach I chose uses Inverse Probability Weighting (IPW) to assign greater weight to control observations that more closely resemble the treated cohort $G_g$. 

Subsequently, to aggregate all $ATT(g,t)$ estimates into a single summary parameter, the authors propose the following aggregation strategy:

\begin{equation} 
\theta = \sum_{g \in \mathcal{G}} \sum_{t=2}^{\mathcal{T}} w(g, t) \cdot ATT(g, t)
\end{equation}

Where $w(g,t)$ represents the weights determined by the researcher. 
In this analysis, I employ a simple averaging scheme mwhere w(g,t) is the same for all treated periods. 
The objective of this comparison is to determine if the Callaway \& Sant’Anna (CS) estimates an Events Studies align with the Two-Way Fixed Effects (TWFE) results. 
A similar behavior between the two models would reinforce the validity of the TWFE estimates and suggest that the possible bias of heterogeneous effects is not a preoccupation \parencite{RiosAvila_Callaway_SantAnna_2021_csdid_slides,Callaway_SantAnna_2021}.

\section{Results}
\label{sec:Results}

\section{Discussion}

\section{Robustness}
spill over effects, placebo tests with other crime types, etc.
Vary the buffers fo the spill over effects, etc.
propensity score matching
treat controls and treated groups as buffers

\section{Conclusions}


\newpage
\printbibliography

\newpage

\section{Appendix}

\appendix
\counterwithin{table}{section}

\section{Detailed description of variables}
\label{sec:variables_descripcion_detallada}

\begin{longtable}{p{0.25\linewidth} p{0.45\linewidth} p{0.25\linewidth}}
    \caption{Data sources}
    \label{tab:variables_descripcion}
    \\
        \toprule
        \textbf{Variable} & \textbf{Description} & \textbf{Data source} \\
        \midrule
        \endhead
        \midrule
        \multicolumn{3}{l}{\textbf{Dependent Variable}} \\
        \midrule
        Proportion of independents & Proportion of independents in the ZAT that year & Movility surveys (2011, 2015, 2019, 2023) \\
        \midrule
        \multicolumn{3}{l}{\textbf{Variable de tratamiento escalonado}} \\
        \midrule
        Tiene OXXO ($=1$) & Si el ZAT tiene OXXO(s) ese año & RUES y Google Maps API \\
        \midrule
        \multicolumn{3}{l}{\textbf{Controles resagados de tiendas de cadena}} \\
        \midrule
        Tiene D1s (=1) & El ZAT tiene D1(s) el periodo anterior al año analizado & RUES y Google Maps API \\
        \midrule
        Tiene ARAs (=1) & El ZAT tiene ARA(s) el periodo anterior al año analizado & RUES y Google Maps API \\
        \midrule
        \multicolumn{3}{l}{\textbf{Controles variables por año}} \\
        \midrule
        Número de personas en el hogar por ZAT & Número promedio de personas en el hogar por ZAT en el año analizado & Encuestas de movilidad (2011, 2015, 2019, 2023) \\
        \midrule
        Ingresos del hogar por ZAT & Ingreso promedio del hogar por ZAT en el año analizado & Encuestas de movilidad (2011, 2015, 2019, 2023) \\
        \midrule
        \multicolumn{3}{l}{\textbf{Controles de línea base}} \\
        \midrule
        Habitantes por localidad en 2007 & Número de habitantes por localidad en el 2007 & Encuesta de Calidad de Vida de Bogotá 2007 \\
        \midrule
        Cantidad de personas por hogar 2007 & Número promedio de personas en el hogar por localidad en el 2007 & Encuesta de Calidad de Vida de Bogotá 2007 \\
        \midrule
        Indice de condiciones de vida 2007 & Indice de condiciones de vida promedio por localidad en el 2007 & Encuesta de Calidad de Vida de Bogotá 2007 \\
        \midrule
        Gasto promedio mensual 2007 & Gasto promedio mensual en el hogar por localidad en el 2007 & Encuesta de Calidad de Vida de Bogotá 2007 \\
        \midrule
        \multicolumn{3}{l}{\textbf{Controles fijos en el tiempo}} \\
        \midrule
        Estrato promedio del ZAT & Estrato promedio del ZAT, que oscila de 1 a 6 & Mapas Bogotá - Infraestructura de Datos Espaciales de Bogotá (IDECA) \\
        \midrule
        Estaciones de transmilenio cercanas & Cantidad de estaciones de transmilenio dentro o cercanos al ZAT & Mapas Bogotá - Infraestructura de Datos Espaciales de Bogotá (IDECA) \\
        \midrule
        Cantidad acceso a vías arteriales & Cantidad de vías arteriales dentro o cercanos al ZAT & Mapas Bogotá - Infraestructura de Datos Espaciales de Bogotá (IDECA) \\
        \midrule
        \multicolumn{3}{l}{\textbf{Control spillover}} \\
        \midrule
        Cantidad de Oxxos cercanos fuera del ZAT & Cantidad de OXXOs fuera del ZAT, que se encuentran cerca de este. & RUES y Google Maps API \\
        \bottomrule
        \multicolumn{3}{p{\linewidth}}{%
        \vspace{0.5em}
        \footnotesize{ \textit{Nota:} No existe controles de línea base por ZAT, por lo que solo se pudo obtener índices previos al 2011 por localidad (ECDV 2007). El control "spillover" tiene en cuenta que hay ZATs que no tienen OXXOs, pero presentan OXXOs cercanos que pueden afectar la dinámica económica y laboral de la zona analizada. Finalmente, cuando de habla de OXXOs, vías arteriales y estaciones de transmilenio cercanos, significa que se encuentran a menos de 400m del ZAT.}}
\end{longtable}

\section{Evolution of OXXO stores in Bogotá divided by Transport Analysis Zones (ZATs)}

\begin{figure}[H]
    \centering
    % Top row
    \begin{subfigure}[b]{0.3\textwidth} % Reduced from 0.4 to 0.3
        \includegraphics[width=\linewidth]{figs_oxxo_maps/mapa_oxxos_binary_zat_2009.png}
        \caption{2009}
        \label{fig:panelA3}
    \end{subfigure}
    \hspace{0.5cm} % Adds specific horizontal space between images
    \begin{subfigure}[b]{0.3\textwidth} 
        \includegraphics[width=\linewidth]{figs_oxxo_maps/mapa_oxxos_binary_zat_2011.png}
        \caption{2011}
        \label{fig:panelB3}
    \end{subfigure}
    \hspace{0.5cm} % Adds specific horizontal space between images
        \begin{subfigure}[b]{0.3\textwidth} 
        \includegraphics[width=\linewidth]{figs_oxxo_maps/mapa_oxxos_binary_zat_2015.png}
        \caption{2015}
        \label{fig:panelC3}
    \end{subfigure}
    
    \vspace{0.3cm} % Vertical space between rows
    
    % Bottom row
    \begin{subfigure}[b]{0.3\textwidth} 
        \includegraphics[width=\linewidth]{figs_oxxo_maps/mapa_oxxos_binary_zat_2019.png}
        \caption{2019}
        \label{fig:panelD3}
    \end{subfigure}
    \hspace{0.5cm}
    \begin{subfigure}[b]{0.3\textwidth} 
        \includegraphics[width=\linewidth]{figs_oxxo_maps/mapa_oxxos_binary_zat_2023.png}
        \caption{2023}
        \label{fig:panelE3}
    \end{subfigure}
    \hspace{0.5cm}
    \begin{subfigure}[b]{0.3\textwidth} 
        \includegraphics[width=\linewidth]{figs_oxxo_maps/mapa_oxxos_binary_zat_2025.png}
        \caption{2025}
        \label{fig:panelF3}
    \end{subfigure}
    
    \caption{\textbf{Evolution of the presence of OXXO by ZAT}}
    \label{fig:fourpanelOXXO}
\end{figure}

\section{Staggered controls of discount store chains}
\label{sec:controles_cadenas_descuento}

\begin{table} [H]
  \centering
  \caption{Means differences of hard discount store staggered controls for the 2018-4 cohort}
  \label{tab:staggered_controls_2018}
  \begin{tabular}{l c c c}
    \toprule
    & \textbf{Still without treatment} & \textbf{Treated} & \textbf{Means Difference} \\
    \midrule
    Has D1s & 0.477 & 0.477 & $0.172^{***}$ \\
    & (0.054) & (0.079) & \\
    \midrule
    Quantity of D1s & 1.233 & 3.250 & $2.017^{***}$ \\
    & (0.214) & (0.543) & \\
    \midrule
    Has Aras & 0.256 & 0.714 & $0.458^{***}$ \\
    & (0.047) & (0.087) & \\
    \midrule
    Quantity of Aras & 0.384 & 1.429 & $0.216^{***}$ \\
    & (0.089) & (0.264) & \\
    \midrule
    Has J\&B & 0.895 & 1.964 & $1.069$ \\
    & (0.050) & (0.096) & \\
    \midrule
    Quantity of J\&B & 0.062 & 0.278 & $0.216^{***}$ \\
    & (0.256) & (1.202) & \\
    \midrule
    UPZs & 86 & 28 & 114 \\
    \bottomrule
  \end{tabular}
  \parbox[t]{\textwidth}{%
    \vspace{0.5em}
    \footnotesize{ \textit{Nota:} Significancia: *** p<0.01, ** p<0.05, * p<0.1.}}
\end{table}


\begin{table} [H]
  \centering
  \caption{Means differences of hard discount store staggered controls for the 2022-4 cohort}
  \label{tab:staggered_controls_2022}
  \begin{tabular}{l c c c}
    \toprule
    & \textbf{Still without treatment} & \textbf{Treated} & \textbf{Means Difference} \\
    \midrule
    Has D1s & 0.491 & 0.881 & $0.390^{***}$ \\
    & (0.067) & (0.042) & \\
    \midrule
    Quantity of D1s & 1.053 & 4.458 & $3.405^{***}$ \\
    & (0.191) & (0.512) & \\
    \midrule
    Has Aras & 0.298 & 0.627 & $0.329^{***}$ \\
    & (0.061) & (0.063) & \\
    \midrule
    Quantity of Aras & 0.386 & 1.712 & $1.326^{***}$ \\
    & (0.089) & (0.287) & \\
    \midrule
    UPZs & 57 & 59 & 116 \\
    \bottomrule
  \end{tabular}
  \parbox[t]{\textwidth}{%
    \vspace{0.5em}
    \footnotesize{ \textit{Nota:} Significancia: *** p<0.01, ** p<0.05, * p<0.1.}}
\end{table}

\end{document}